\subsection{A sharp query law for a cache with two keys}
\label{sec:rankonequerylaw}

A fast cache index can remove the length of the cached prefix from
the query cost. It cannot generally remove the width of a fresh
query. The following family makes this distinction quantitative.

\begin{theorem}[Two-key attention]
\label{thm:rankonequerylaw}
Let $n\ge1$, and let $\beta\ge0$ be dyadic or equal to $\sqrt n$.
A cache contains equal numbers of the
keys $\mathbf1$ and $-\mathbf1$, with respective scalar values $+1$
and $-1$. The fresh query is $x\in\{-1,+1\}^n$, accessed by coordinate
probes. Scores are $\beta\langle x,k\rangle/n$. Apply one tanh row
to the attention value and use the two logits given by that row and
its negative. Let $Q^*(n,\beta)$ be the infimum, over exact samplers,
of the worst-case expected number of fresh-query probes. Arbitrary
preprocessing independent of $x$ is allowed. Then, with absolute
implicit constants,
\[
 Q^*(n,\beta)=\Theta\!\left(\min\{n,\beta+\beta^2\}\right).
\]
There is a finite-bit implementation with expected online work
\[
 O\!\left(B^4\left[1+\min\{n,\beta+\beta^2\}\right]\right),
\]
where $B$ contains the parameter encoding lengths and $\log(n+2)$.
The cache can be stored by its two key types and multiplicities.
At standard scaling, the fresh-query probe complexity is $\Theta(n)$,
independently of the number of cached copies.
\end{theorem}
\paragraph{Proof idea.}
The two-key cache turns the query mean into a scalar tanh composition.  Factories give the upper bound.  For the lower bound, a product-distributed query and the transcript score relate the output derivative to the number of coordinates read.

\begin{proof}
Write $M=n^{-1}\sum_i x_i$ and $F_3=\tanh\circ\tanh\circ\tanh$.
The attention value is $\tanh(\beta M)$, so the final token has
mean $F_3(\beta M)$ when written as a sign.

For the upper bound, obtain a source sign of mean $M$ by querying
one uniform coordinate. The arbitrary-radius tanh factory uses at
most $2\max\{\beta,\beta^2\}$ such requests, and each of the next
two unit-radius tanh factories uses at most two requests to its
predecessor. Fresh randomness gives the factor eight. Repeated
input coordinates may be cached, which can only reduce the number
of probes. Alternatively read all $n$ input coordinates and make a
certified scalar comparison. Choose the cheaper construction from
the known parameters. The quadratic-radius bit bound and the
precision stopping lemma give the stated finite-bit work.

For $0\le\beta\le1$, the output laws on $x=\mathbf1$ and
$x=-\mathbf1$ have total variation distance $F_3(\beta)$.
Couple all preprocessing and online random bits. Until the first
probe the executions are identical, so the probability of making
at least one probe is at least that distance. The inequalities
$\tanh u\ge u/2$ for $u\in[0,1]$ give
$F_3(\beta)\ge\beta/8$. This proves the required bound in this range,
since $\beta+\beta^2\le2\beta$.

For $\beta\ge1$, draw the query coordinates independently with
common mean $t$, and let $H(t)$ be the mean of the output sign.
At $t=0$, put $S=\sum_i x_i$ and $Z=S/\sqrt n$. Differentiating
the finite product distribution gives
\[
 H'(0)=\mathbb E_0[F_3(\beta S/n)S]
 =\sqrt n\,\mathbb E_0\left[Z F_3\left(\frac\beta{\sqrt n}Z\right)\right].
\]
Here $\mathbb E Z^2=1$ and $\mathbb E Z^4=3-2/n\le3$.
Paley--Zygmund's inequality, or its elementary Cauchy--Schwarz proof,
therefore gives
\[
 \Pr(|Z|\ge1/2)\ge\frac{(1-1/4)^2}{3}=\frac3{16}.
\]
Let $d=\min\{\beta/\sqrt n,1\}$. On this event, oddness and
monotonicity imply
\[
 ZF_3\left(\frac\beta{\sqrt n}Z\right)
 \ge\frac12 F_3(d/2)\ge\frac d{32}.
\]
The integrand is nonnegative everywhere, hence
$H'(0)\ge3\min\{\beta,\sqrt n\}/512$.
Lemma~\ref{lem:transcript} and Cauchy--Schwarz give
\[
 \sup_x\mathbb E_xQ\ge\mathbb E_0Q
 \ge H'(0)^2
 \ge\frac1{32768}\min\{\beta^2,n\}.
\]
Finally $\beta+\beta^2\le2\beta^2$ in this range. Combining both
ranges gives the universal lower constant $1/65536$.
The case $\beta=0$
is a fair output coin and uses no input probes.
\end{proof}

\begin{theorem}[Worst-case query cost for a rank-one cache]
\label{thm:rankone-worstlaw}
Let $T,n\ge2$, and let $\beta\ge1$ be dyadic or equal to $\sqrt n$.
Consider all length-$T$ caches with finitely encoded dyadic keys and
values, whose keys
lie in $[-1,1]^n$ on one affine line, whose scalar values lie in
$[-1,1]$, and whose fresh query is $x\in\{-1,1\}^n$ accessed by
coordinate probes. The required output sign has mean
$\tanh(\tanh Y(x))$, where $Y$ is the attention-coordinate mean
with scores $\beta\langle x,k_j\rangle/n$.
Allow preprocessing independent of $x$ and polynomial in the cache
length, width, and parameter encoding size.

The optimal worst-cache, worst-query expected probe cost is
\[
 \Theta\!\left(\min\{n,\beta^2\}\right),
\]
with absolute implicit constants.  The upper bound has expected online bit work
\[
 O\!\left([1+\min\{n,\beta^2\}+\log T]B^4\right)
\]
after polynomial $(T,n,B)$ prefill, where $B\ge16$ includes all parameter
encoding lengths and $\log_2(T+n+2)$. At standard scaling the sharp
fresh-query probe order is $n$, independently of the prefix length.
The theorem is a worst-case law over the stated class of caches.
\end{theorem}
\paragraph{Proof idea.}
The two-key construction supplies the lower bound for every cache length.  For the upper bound, express an arbitrary line cache using its two extreme keys and expose the resulting one-dimensional query through source signs; use the explicit index when reading the whole query is cheaper.

\begin{proof}
For even $T$, the equal-multiplicity two-key cache in
Theorem~\ref{thm:rankonequerylaw} supplies the lower bound even
with arbitrary preprocessing independent of $x$.
For odd $T$, use $T-1$ equally split keys and one zero key with
value zero. Its attention output has the same sign as
$\tanh(\beta M)$ and magnitude at least
$(2/3)|\tanh(\beta M)|$. The map
$\phi(u)=\tanh(\tanh u)$ is increasing and concave for $u\ge0$,
so the absolute output mean is at least $(2/3)|F_3(\beta M)|$.
The first-score calculation in that theorem therefore gives
$H'(0)\ge\min\{\beta,\sqrt n\}/256$.
Lemma~\ref{lem:transcript} proves the same lower constant
$1/65536$ in this case as well.

For the upper bound, find the two extreme keys along the affine line.
Their midpoint is $c$ and their half-difference is $u$. Then
\[
 k_j=c+t_ju,\qquad |t_j|\le1,\qquad \|u\|_\infty\le1.
\]
If all keys agree, the attention value is independent of the query.
Otherwise each $t_j$ is a rational ratio of two dyadic coordinate
differences, so its bit length is $O(B)$. A fresh sign of mean
$z=\langle x,u\rangle/n$ needs at most one input probe: select
coordinate $i$ with mass $|u_i|/n$, return its signed input, and
complete the unused mass with a fair sign. Independent selections
give independent source signs for every fixed query.

The rank-one reduction after Theorem~\ref{thm:implicitrankone}, with
$R=n$ and $p=(1+z)/2$, gives slopes $a_j=2\beta t_j$ and a rational
slope bound $2\beta\le\Lambda\le4\beta$.
Its radius-two attention sign needs fewer than
$45\Lambda^2\le720\beta^2$ source requests.
The arbitrary-radius tanh factory in
Theorem~\ref{thm:quadraticradiusbits} needs at most eight such signs
to generate a sign of mean $\tanh Y$, and at most two further
calls produce a sign of mean $\tanh(\tanh Y)$.
Predictable charging gives at most $11520\beta^2$ input requests.
Its expected bit work is $O(\beta^2B^4)$.

Alternatively, read all $n$ input signs and use the explicit
rank-one index in Theorem~\ref{thm:newrankone}. Its work is
$O((n+\log T)B^4)$ and its probe count is $n$.
Use the implicit construction only when $12000\beta^2<n$.
Then its prefill factor $\Lambda^4$ is $O(n^2)$, so its
$O(T^3(1+\Lambda^4)B^4)$ preprocessing is polynomial in
$T,n,B$. In the other case the near-linear explicit index
suffices. Thus no preprocessing exponential in the encoding of
a numerically huge temperature is being claimed.
This proves the two upper bounds and almost-sure termination.
\end{proof}

\paragraph{Adding the current position.}
The affine-line index also tolerates a fixed number of explicitly
available exceptional keys, such as the current position's own key.
Use the monotone geometric blocks for the line and one singleton
block per exception. Shift by the largest weight among the line and
the exceptions. The ideal envelope is at most twice the total mass;
the same dyadic floor argument gives at most $17/8$ proposals.
The additional query work is proportional to the number of exceptional
inner products. Thus including the current position does not require
its key to lie on the cached affine line.
